% main.tex -- Letter (PRL style), GYNI causal inequality
% Companion files: supplement.tex (Supplemental Material), refs.bib.
% All decimals are directed-rounded from the stored exact rationals:
% lower bounds rounded DOWN, upper bounds rounded UP (see Supplemental Material, Sec. S8).
\documentclass[aps,prl,twocolumn,superscriptaddress,longbibliography,nofootinbib]{revtex4-2}

\usepackage{amsmath,amssymb,amsthm}
\usepackage{bm}
\usepackage[colorlinks=true,linkcolor=blue,citecolor=blue,urlcolor=blue]{hyperref}

\newtheorem{theorem}{Theorem}
\newtheorem{lemma}{Lemma}

\newcommand{\Tr}{\operatorname{Tr}}
\newcommand{\ket}[1]{\lvert #1\rangle}
\newcommand{\bra}[1]{\langle #1\rvert}
\newcommand{\one}{\mathbb{I}}
\newcommand{\tp}{\mathsf{T}}
\newcommand{\CC}{\mathbb{C}}
\newcommand{\RR}{\mathbb{R}}
\newcommand{\ZZ}{\mathbb{Z}}
\newcommand{\Igy}{I_{\mathrm{GYNI}}}
\newcommand{\Imax}{I_{\mathrm{GYNI}}^{\max}}
\newcommand{\cW}{\mathcal{W}}

\begin{document}

\title{Pinning down the maximal quantum violation of the Guess-Your-Neighbour's-Input causal inequality}

\author{Ansh Mishra}
\author{Aryan Senthilkumar}
\noaffiliation

\date{\today}

\begin{abstract}
Causal inequalities bound the correlations that parties acting in a definite causal order can produce, and process matrices---quantum theory without a global causal order---can violate them. Determining the maximal violations, the analogues of Tsirelson's bound, is open in general. For the Guess-Your-Neighbour's-Input (GYNI) inequality, whose causal bound is $1/2$, the maximal success probability over finite-dimensional process-matrix strategies was known only to lie between the numerical see-saw value $0.6219$ and the upper bound $0.7592$. We prove $0.622165901353\le\Imax\le0.622183755532$, narrowing this interval by a factor of about $7700$. The upper bound comes from a new hierarchy of semidefinite relaxations, built on a L\"uders normal form for local instruments and on a moment matrix indexed by words in the infinite dihedral group; at level 8 it is certified by an exact rational dual certificate checked in integer arithmetic. The lower bound is an explicit process with four coherently superposed qubit Jordan blocks, verified in exact rational arithmetic. Both bounds, including the validity of the relaxation, are machine-checked in the Lean~4 proof assistant.
\end{abstract}

\maketitle

\emph{Introduction.}---Correlations between parties whose operations are embedded in a definite causal order, or in a probabilistic mixture of such orders, satisfy \emph{causal inequalities}~\cite{Oreshkov2012,Branciard2016,OreshkovGiarmatzi2016,Abbott2016}. Oreshkov, Costa and Brukner showed that the process-matrix framework, which assumes quantum theory in each local laboratory but no global causal structure, predicts violations of such inequalities~\cite{Oreshkov2012}; see Ref.~\cite{Costa2026review} for a review, and Ref.~\cite{Chiribella2013} for a related model of indefinite causal order, the quantum switch. How large can these violations be? For Bell inequalities the analogous question is answered by Tsirelson bounds~\cite{Tsirelson1980}, and the NPA hierarchy of semidefinite programs (SDPs)~\cite{NPA2007,NPA2008} computes them systematically. For causal inequalities, upper bounds on process-matrix violations have been derived only recently~\cite{Brukner2015,LiuChiribella2025}. Liu and Chiribella derived an SDP bound for every causal inequality, determined the exact maximal violations for a class that includes a biased version of the inequality of Ref.~\cite{Oreshkov2012} and the ``lazy'' GYNI inequality (LGYNI)~\cite{Branciard2016}, and suggested that their relaxation may be the first level of a hierarchy~\cite{LiuChiribella2025}.

The Guess-Your-Neighbour's-Input (GYNI) inequality and LGYNI represent the two families of nonequivalent causal inequalities of the simplest bipartite scenario, with binary inputs and outputs~\cite{Branciard2016}. In the GYNI game Alice and Bob receive independent, uniformly random bits $x$ and $y$ and output bits $a$ and $b$; they win if each guesses the other's input:
\begin{equation}
\Igy=\frac14\sum_{x,y\in\{0,1\}}p(a=y,b=x\,|\,x,y)\le\frac12
\label{eq:gyni}
\end{equation}
for all causal correlations~\cite{Branciard2016}. For the maximal value $\Imax$ over process matrices, the best upper bound valid in all dimensions was $0.7592$~\cite{LiuChiribella2025}. See-saw optimisation in fixed local dimension $d=d_{A_I}=d_{A_O}=d_{B_I}=d_{B_O}$ found $0.6218$ for $d=5$~\cite{Branciard2016} and $0.6219$ for $d=6,7$ (in floating-point arithmetic), with no significant gain up to $d=8$~\cite{BoghiuSimonov2026}. The value of $\Imax$ was open~\cite{Branciard2016,LiuChiribella2025,BoghiuSimonov2026}.

In this Letter we prove
\begin{equation}
0.622165901353\;\le\;\Imax\;\le\;0.622183755532 ,
\label{eq:main}
\end{equation}
where $\Imax$ is the supremum over strategies whose local systems have finite, but arbitrary, dimension; local ancillas and shared entanglement are included. The interval has width $1.79\times10^{-5}$, compared with $0.137$ before. The upper bound comes from a new NPA-type hierarchy for process-matrix correlations, whose validity we prove (Theorem~\ref{thm:sound}), evaluated at level~8. The lower bound is an explicit strategy in which the two measurements of each party are built from four coherently superposed qubit ``Jordan blocks''. Both bounds are certified by exact computations, which stand-alone programs check in integer or rational arithmetic without any numerical solver~\cite{repo}.

\emph{Setting.}---Alice has input and output spaces $A_I,A_O$, Bob has $B_I,B_O$, all of finite dimension. An instrument $\{M_{a|x}\}_a$ is a family of completely positive (CP) maps from operators on $A_I$ to operators on $A_O$ such that $\sum_aM_{a|x}$ is trace preserving (TP); likewise $\{N_{b|y}\}_b$ for Bob. We use Choi operators $C_F=\sum_{ij}\ket{i}\bra{j}\otimes F(\ket{i}\bra{j})$~\cite{Choi1975}. A \emph{process} is an operator $W\succeq0$ on $A_I\otimes A_O\otimes B_I\otimes B_O$ with $\Tr[W(C_A\otimes C_B)]=1$ for all Choi operators $C_A,C_B$ of CPTP maps~\cite{Oreshkov2012,Araujo2015}; equivalently, $W\succeq0$, $\Tr W=d_{A_O}d_{B_O}$, and $W$ lies in the subspace of valid processes~\cite{Araujo2015}. A strategy $S=(W;M;N)$ yields $p(a,b|x,y)=\Tr[W(C_{M_{a|x}}\otimes C_{N_{b|y}})]$; we write $\Igy(S)$ for the left-hand side of Eq.~\eqref{eq:gyni} and $\Imax=\sup_S\Igy(S)$. Local ancillas and shared entangled states are covered, as they can be absorbed into $W$. (Our Choi convention differs from that of Ref.~\cite{Oreshkov2012} by the transposition $W\mapsto W^\tp$, which maps processes onto processes and leaves $\Imax$ unchanged.)

\emph{L\"uders normal form.}---The first step removes the generality of the instruments.

\begin{lemma}[L\"uders normal form]\label{lem:lueders}
For every strategy there are finite-dimensional spaces $H_A,H_B$, projective measurements $\{P_{a|x}\}_a$ on $H_A$ and $\{Q_{b|y}\}_b$ on $H_B$, and a process $\widetilde W$ with inputs $H_A,H_B$ and outputs $H_A\otimes X$, $H_B\otimes Y$ ($X=Y=\CC^2$), which reproduces $p(a,b|x,y)$ with the L\"uders instruments
\begin{equation}
\widetilde M_{a|x}(\rho)=P_{a|x}\rho P_{a|x}\otimes\ket{x}\bra{x}_X ,\quad
\widetilde N_{b|y}(\sigma)=Q_{b|y}\sigma Q_{b|y}\otimes\ket{y}\bra{y}_Y .
\label{eq:lueders}
\end{equation}
Moreover, $\widetilde W$ is block diagonal in the register values.
\end{lemma}

The proof [Supplemental Material (SM), Sec.~S3] dilates each instrument as $M_{a|x}(\rho)=\sum_eT_{x,e}P_{a|x}J\rho J^\dagger P_{a|x}T_{x,e}^\dagger$, with a fixed isometry $J$ from $A_I$ into $H_A$, projectors $P_{a|x}$ obtained from a unitary completion of the Stinespring isometry of $\{M_{a|x}\}_a$, and Kraus operators $T_{x,e}$ of a channel. The $x$-independent pre-processing $J$ and the post-processing, which is controlled by the setting register, are then absorbed into $\widetilde W=(\Phi_A^\dagger\otimes\Phi_B^\dagger)(W)$, where $\Phi_A(C)=\sum_{x,e}\hat K_{x,e}C\hat K_{x,e}^\dagger$ with $\hat K_{x,e}=J^\tp\otimes T_{x,e}\otimes\bra{x}_X$ maps Choi operators of channels to Choi operators of channels. The transpose is essential: pre-composition with $J$ acts on Choi operators as $C\mapsto(J^\tp\otimes\one)C(J^\tp\otimes\one)^\dagger$, while $J^\dagger$ in place of $J^\tp$ would implement pre-composition with the complex-conjugate isometry $\bar J$, which for complex $J$ does not reproduce the statistics.

\emph{Moment hierarchy.}---In L\"uders form the measurements of each party are described by two unitary involutions $O_x=P_{0|x}-P_{1|x}$. By the universal property of the free product, $g_x\mapsto O_x$ defines a unitary representation $\pi_A$ of the infinite dihedral group $G=\ZZ_2*\ZZ_2=\langle g_0,g_1\,|\,g_0^2=g_1^2=e\rangle$, whose elements are the reduced words $w=(x_1,\dots,x_n)$ with $x_i\neq x_{i+1}$. Let $\cW_L$ be the $2L+1$ words of length at most $L$. For $k=(w,r)\in\cW_L\times\{0,1\}$ we define the word vector
\begin{equation}
\ket{k}=(\one\otimes\pi_A(w))\ket{\Phi}\otimes\ket{r}\in H_A\otimes H_A\otimes X ,
\end{equation}
with $\ket{\Phi}=\sum_i\ket{i}\ket{i}$, and analogously $\ket{l}$ for $l=(v,s)$ on Bob's side. The level-$L$ moment matrix is the Gram matrix of $\widetilde W$ on these vectors,
\begin{equation}
\Gamma[(k,l),(k',l')]=\bra{k,l}\widetilde W\ket{k',l'} .
\label{eq:gamma}
\end{equation}
The operator $\ket{k'}\bra{k}$ is the Choi operator of $\rho\mapsto\pi_A(w')\rho\,\pi_A(w)^\dagger\otimes\ket{r'}\bra{r}$, and its partial trace over the output equals $\delta_{rr'}[\pi_A(w^{-1}w')]^\tp$. Hence a combination $C(c)=\sum c_{kk'}\ket{k'}\bra{k}$ over pairs with $r=r'$ satisfies $\Tr_{\rm out}C(c)=\lambda\one$ \emph{in every realisation} whenever $\sum c_{kk'}\,w^{-1}w'=\lambda e$ holds in the group algebra $\RR[G]$; we then say that $c$ has scalar trace $\lambda$. Validity of the process enters through the following lemma.

\begin{lemma}[Scalar-trace factorisation]\label{lem:v2}
Let $W$ be a process and let $C_A$, $C_B$ be arbitrary, not necessarily Hermitian, operators with $\Tr_{A_O}C_A=\lambda_A\one$ and $\Tr_{B_O}C_B=\lambda_B\one$. Then $\Tr[W(C_A\otimes C_B)]=\lambda_A\lambda_B$.
\end{lemma}

Indeed, every such $C_A$ is a complex linear combination $\sum_i\alpha_iC_i$ of Choi operators of CPTP maps with $\sum_i\alpha_i=\lambda_A$ (SM, Sec.~S3), and the claim follows by bilinearity.

The certified level-$L$ relaxation consists of the following constraints on a real symmetric $\Gamma$, with $\langle c\otimes d,\Gamma\rangle:=\sum c_{kk'}d_{ll'}\Gamma[(k,l),(k',l')]$:
(V1)~$\Gamma\succeq0$. Since $\Gamma$ vanishes unless $r=r'$ and $s=s'$ (Lemma~\ref{lem:lueders}), this amounts to four positive semidefinite (PSD) real blocks $\Gamma_{rs}$ of size $n_L=(2L+1)^2$.
(V2)~$\langle c\otimes d,\Gamma\rangle=\lambda_A\lambda_B$ whenever $c$ and $d$ have scalar traces $\lambda_A$ and $\lambda_B$. These constraints are spanned by differences of pairs with the same group element $w^{-1}w'$ (``trace annihilating'') combined with such differences or with the pair of empty words (``trace preserving'') on the other side.
(S)~Invariance under the symmetries of the game: party exchange and the relabellings $(x\to1-x,\,b\to1-b)$ and $(y\to1-y,\,a\to1-a)$. In the dihedral word basis they act on $\Gamma$ by signed permutations; e.g., the first exchanges the letters of Alice's words, flips her register and multiplies Bob's word vectors by $(-1)^{|v|}$.
Finally, $P_{a|x}=[\pi_A(e)+(-1)^a\pi_A(g_x)]/2$, so every probability, and hence $\Igy$, is a linear function $o(\Gamma)$.

\begin{theorem}[Soundness]\label{thm:sound}
For every $L\ge1$ and every strategy $S$ with finite-dimensional local systems there is a real $\Gamma$ satisfying (V1), (V2) and (S) with $o(\Gamma)=\Igy(S)$. Hence $\Imax\le\beta_L:=\sup\{o(\Gamma):\text{(V1), (V2), (S)}\}$.
\end{theorem}

The proof (SM, Sec.~S3) combines Lemmas~\ref{lem:lueders} and \ref{lem:v2} with two further observations: the real part of $\Gamma$ is again feasible, and averaging over the finite group generated by the three symmetries preserves feasibility and the objective. Only the relations $g_x^2=e$ are used, never a relation specific to particular projectors, so the constraints hold in every dimension. The constraints of level $L$ are contained in those of level $L+1$, so $\beta_L$ does not increase with $L$.

The positivity (V1) is the key new ingredient. Numerically, level~1 already gives $\beta_1\approx0.7463<0.7592$, while at level~2 dropping (V1) but imposing validity of the canonical processes of Ref.~\cite{LiuChiribella2025}, which are linear images of $\Gamma$, reproduces $0.7592$. The certified program does not contain these canonical-process constraints (at level~2, adding them does not change the numerical value), and none of our bounds relies on them.

\emph{Exact certificates.}---We solve the relaxations numerically (Table~\ref{tab:levels}) and then certify each bound with an exact dual certificate. Write the linear constraints (V2) and (S) as $Az=b$, where $z$ collects the independent entries of the four blocks $\Gamma_k(z)$, and $o(\Gamma)=o\cdot z$. For multipliers $\lambda$, define symmetric matrices $Y_k(\lambda)$ by stationarity, $\sum_k\langle Y_k,\Gamma_k(z)\rangle=(A^\tp\lambda-o)\cdot z$ for all $z$. Every feasible $z$ then satisfies
\begin{equation}
o\cdot z=\lambda\cdot b-\sum_k\langle Y_k,\Gamma_k(z)\rangle\le\lambda\cdot b\quad\text{if all }Y_k\succeq0 .
\label{eq:weakdual}
\end{equation}
We shift a numerical dual solution by $\varepsilon\one$ to make it strictly feasible ($\varepsilon=10^{-7}$ for $5\le L\le7$ and $\varepsilon=10^{-9}$ at $L=8$), solve for multipliers of all constraints by least squares, and round them to $\lambda\in2^{-50}\ZZ^m$. The $Y_k$ are then \emph{defined} exactly from $\lambda$, and their positive definiteness is checked with Sylvester's criterion: all leading principal minors of the integer matrices $2^{54}Y_k$ are computed exactly by fraction-free Bareiss elimination~\cite{Bareiss1968} and found to be positive. The bound $\lambda\cdot b$ is an exact rational number. (Rounding numerical solutions to exact certificates is standard in sums-of-squares optimisation~\cite{PeyrlParrilo2008}.) Every diagonal entry of $\Gamma$ belongs to a trace-preserving pair and is fixed to $1$ by (V2), so $\Tr\Gamma_k=n_L$ on the feasible set; the margin therefore raises the bound by $\varepsilon n_L$ relative to the numerical dual value.

At level~8 the certified program has $167{,}620$ variables, $m=789{,}678$ linear constraints and four blocks of size $n_8=289$. The certificate gives
\begin{equation}
\beta_8\le\frac{175129158097837}{2^{48}}=0.6221837555310\ldots ,
\end{equation}
which proves the upper bound in Eq.~\eqref{eq:main}; the margin costs only $\varepsilon n_8=2.89\times10^{-7}$. A stand-alone verifier rebuilds the constraints, recomputes the $Y_k$ from the stored integers and checks all minors in about $8$ minutes on one core; the decisive steps use exact integer arithmetic only~\cite{repo}. Beyond the proof, we evaluated all $789{,}678$ level-8 constraints on the moment matrices of more than $100$ strategies---general instruments passed through a separately written implementation of Lemma~\ref{lem:lueders}, generic L\"uders-form processes, see-saw optima, and the strategies below---and found no residual above $2.1\times10^{-14}$ (SM, Sec.~S6).

\begin{table}[t]
\caption{Certified upper bounds on the level-$L$ value $\beta_L$, and hence on $\Imax$ (rounded up). $n_L=(2L+1)^2$ is the size of each of the four PSD blocks and $m_L$ the number of linear constraints. The last column is the optimum returned by an interior-point solver (numerical, not rigorous); the certified bounds exceed it essentially by the margin $\varepsilon n_L$, with $\varepsilon=10^{-6}$ at $L=4$, $10^{-7}$ at $L=5,6,7$ and $10^{-9}$ at $L=8$. Entries marked $\dagger$ were certified in the equivalent formulation with words in the two projectors $P_{0|x}$ (SM, Sec.~S4).}
\label{tab:levels}
\begin{ruledtabular}
\begin{tabular}{rrrll}
$L$ & $n_L$ & $m_L$ & certified bound & numerical\\
\hline
1 & 9   & --                    & --                          & 0.7462629\\
2 & 25  & $5{,}486^\dagger$     & $0.6434871^\dagger$         & 0.6434843\\
3 & 49  & $21{,}708^\dagger$    & $0.6267932^\dagger$         & 0.6267499\\
4 & 81  & $59{,}326$            & $0.6233558$                 & 0.6232747\\
5 & 121 & $134{,}604$           & $0.6224213$                 & 0.6224092\\
6 & 169 & $265{,}742$           & $0.6222569$                 & 0.6222399\\
7 & 225 & $475{,}300$           & $0.6222199$                 & 0.6221974\\
8 & 289 & $789{,}678$           & $\mathbf{0.622183755532}$   & 0.6221835\\
\end{tabular}
\end{ruledtabular}
\end{table}

\emph{Explicit strategy.}---By Jordan's lemma~\cite{Pironio2009}, two binary projective measurements decompose the local space into invariant blocks of dimension at most two, each characterised by one angle. This motivates the following family of L\"uders-form strategies, identical for Alice and Bob: $H=\CC^2\otimes\CC^J$ (a ``Jordan qubit'' and a label), $A_I=H$, $A_O=H\otimes\CC^2$, the instruments~\eqref{eq:lueders}, and
\begin{equation}
P_{0|x}=\sum_{j=1}^J\ket{\phi_j^x}\bra{\phi_j^x}\otimes\ket{j}\bra{j},\;\;
\ket{\phi_j^x}=\cos\tfrac{t_j}{2}\ket{0}+(-1)^{x+1}\sin\tfrac{t_j}{2}\ket{1},
\label{eq:jordan}
\end{equation}
so that the two measurement vectors in block $j$ enclose the angle $t_j$, $\langle\phi_j^0|\phi_j^1\rangle=\cos t_j$. For fixed angles the optimal process is the solution of an SDP, which we reduce with symmetries (twirls of the register and of the label phases, and the GYNI group); we optimise the angles with gradients obtained from the envelope theorem.

A single block ($J=1$) reaches only $\approx0.6067$ numerically; the higher values in Table~\ref{tab:strategies} come from processes that are coherent across blocks with different angles. For $J=4$ the angles are $(t_1,\dots,t_4)=(9.59850^\circ,34.42477^\circ,51.18469^\circ,84.13806^\circ)$, chosen with rational $\tan(t_j/4)$ so that all entries of the $P_{a|x}$ are rational. The process is
\begin{equation}
W=\frac{1}{64}\one+(1-2^{-30})\sum_{t=1}^{1895}\frac{y_t}{2^{40}}E_t ,
\label{eq:Wexplicit}
\end{equation}
with integers $y_t$ and integer matrices $E_t$; each $E_t$ is a symmetrised tensor product of two local operators, each of which lies in a single Hilbert--Schmidt sector of $A_I\otimes A_O$ (resp.\ $B_I\otimes B_O$). The factor $1-2^{-30}$ mixes in a small amount of the maximally mixed process to make $W$ positive definite. $W$ is a real $16384\times16384$ matrix with $\Tr W=256=d_{A_O}d_{B_O}$. A stand-alone verifier checks in exact arithmetic that (i) every $E_t$ lies in the subspace of valid processes~\cite{Oreshkov2012,Araujo2015}, (ii) $W$ is block diagonal with $676$ blocks, each positive definite (Bareiss), (iii) the instruments are valid, and (iv) the value is
\begin{equation}
\Igy=0.62216590135390844\ldots ,
\end{equation}
a rational number with a $548$-bit denominator. This proves the lower bound in Eq.~\eqref{eq:main}. The check takes under $40$ minutes on one core, uses no symmetry assumption, and is independent of the hierarchy.

\begin{table}[t]
\caption{Explicit $J$-block strategies of Eq.~\eqref{eq:jordan}, with $d_{A_I}=2J$ and $d_{A_O}=4J$ (including the setting register); angles rounded to $0.1^\circ$, values rounded down. v: exact value, and all checks of the stand-alone verifier passed without symmetry assumptions. e: exact value, obtained with an earlier, symmetry-reduced positivity check (SM, Sec.~S5). n: numerical optimum.}
\label{tab:strategies}
\begin{ruledtabular}
\begin{tabular}{cll}
$J$ & angles $t_j$ (degrees) & $\Igy$\\
\hline
1 & 36.1                   & $0.6067^{\rm n}$\\
2 & 22.7, 70.3             & $0.621789802^{\rm e}$\\
3 & 13.8, 40.2, 82.7       & $0.622146712690^{\rm v}$\\
4 & 9.6, 34.4, 51.2, 84.1  & $\mathbf{0.622165901353}^{\rm v}$\\
\end{tabular}
\end{ruledtabular}
\end{table}

\emph{Discussion.}---Equation~\eqref{eq:main} determines $\Imax$ to within $1.79\times10^{-5}$. Nearly all of this width is the gap of $1.76\times10^{-5}$ between the numerical level-8 optimum, $0.6221835$, and our lower bound; the margin of the level-8 certificate accounts for only $2.89\times10^{-7}$. No certificate for level~8 can go below the level-8 value $\beta_8$, so a narrower interval requires higher levels ($L\ge9$; the level-8 solve took $220$~s with a custom interior-point method) or larger $J$. Whether $\beta_L$ converges to $\Imax$ is open: unlike for the NPA hierarchy~\cite{NPA2008}, we have no convergence theorem, and the moment matrices we computed are not flat (e.g., numerical rank $134$ of $169$ at $L=6$, versus $109$ of $121$ on the words of length at most $5$), so no strategy can be read off from them. A geometric (Aitken) extrapolation of the numerical values in Table~\ref{tab:levels} gives about $0.62218$; such extrapolations are heuristic (other fits we tried are already excluded by our lower bound), and we make no claim about the exact value of $\Imax$.

The method is not specific to GYNI. Lemmas~\ref{lem:lueders} and~\ref{lem:v2} do not refer to the game, and our implementation handles any number of settings with binary outcomes. As consistency checks, the relaxation reproduces numerically, to solver precision, the exact LGYNI value $0.8194$ of Ref.~\cite{LiuChiribella2025} at level~2 and, with a suitable word set, the maximal success probability $(2+\sqrt2)/4$ of the game of Ref.~\cite{Oreshkov2012}, proven optimal in Ref.~\cite{LiuChiribella2025}. We expect that Lemma~\ref{lem:lueders}, applied party by party, together with a multilinear version of Lemma~\ref{lem:v2}, extends the hierarchy to multipartite causal inequalities~\cite{BaumelerWolf2014,Abbott2016}, where the structure of valid processes is richer~\cite{Araujo2015,OreshkovGiarmatzi2016}; working this out, and comparing with block-matrix relaxations~\cite{DAlessandro2026}, is left for future work. Our bounds concern processes of finite dimension; Lemma~\ref{lem:lueders} does not cover infinite-dimensional processes.

\emph{Formal verification.}---Both bounds in Eq.~\eqref{eq:main} are machine-checked in Lean~4 with Mathlib~\cite{Lean4}, for every strategy with finite-dimensional local systems, using the standard definitions of processes and instruments.
\begin{itemize}
\item \emph{Lower bound.} The theorem \texttt{gyni\_lower\_bound\_exact} exhibits the explicit process and instruments, proves that they are valid, and proves that their GYNI value equals the rational number above. Positive semidefiniteness is established by kernel-checked integer certificates together with a proved party-swap symmetry.
\item \emph{Upper bound.} The theorem \texttt{gyni\_upper\_bound} states $\Igy\le175129158097837/2^{48}$. Its proof formalises the whole argument: Lemma~\ref{lem:lueders} (with an explicit Stinespring dilation and unitary completion), Lemma~\ref{lem:v2}, the constraints (V1) and (V2) on the moment matrix, weak duality, and the level-8 certificate, whose positivity is checked by the Lean kernel. Before the formal check, the dual certificate is averaged over the symmetry group of the game. The averaged certificate is valid for the relaxation without the constraints (S), with the same bound, so the symmetry reduction is not needed in the formal proof.
\end{itemize}
All theorems depend only on Lean's standard axioms.

All code, certificates, verification logs and tests are available in Ref.~\cite{repo}.

\bibliography{refs}

\end{document}
